\section*{Chapter \ref{ch:b2:alkene-redox} --- Oxidation and Reduction of Alkenes}

\begin{solution}{exo:b2:alkene-redox:1}
Methylcyclohexane; 1-methylcyclohexan-1-ol; trans-2-methylcyclohexan-1-ol (syn addition of
H and B, so the \ce{OH} and the new H on the same side, the methyl trans to the \ce{OH});
1-methylcyclohexene oxide; cis-1-methylcyclohexane-1,2-diol; 6-oxoheptanal
\ce{CH3CO(CH2)4CHO}.
\end{solution}

\begin{solution}{exo:b2:alkene-redox:2}
Hydroboration--oxidation: 2-methylpropan-1-ol. Acid-catalysed hydration:
2-methylpropan-2-ol.
\end{solution}

\begin{solution}{exo:b2:alkene-redox:3}
trans-2,3-Dimethyloxirane, the two methyls on opposite sides of the ring. It is chiral (no
mirror plane) and is formed as a racemic mixture.
\end{solution}

\begin{solution}{exo:b2:alkene-redox:4}
Reducing work-up: two molecules of propanal. Oxidising work-up: two molecules of propanoic
acid.
\end{solution}

\begin{solution}{exo:b2:alkene-redox:5}
(a) meso-butane-2,3-diol, inactive because it is achiral. (b) the
(2\textit{R},3\textit{R}) and (2\textit{S},3\textit{S}) diols in equal amounts, inactive
because racemic.
\end{solution}

\begin{solution}{exo:b2:alkene-redox:6}
In acid water attacks the tertiary carbon, in base hydroxide attacks the \ce{CH2}; in both
cases an \ce{OH} ends on each carbon: the same 2-methylpropane-1,2-diol. With methanol the two
routes give different ethers (as in the chapter).
\end{solution}

\begin{solution}{exo:b2:alkene-redox:7}
Hydrogenation with one equivalent of \ce{H2} over a poisoned palladium catalyst
(Lindlar): syn addition to the triple bond stops at the (\textit{Z})-alkene.
\end{solution}

\begin{solution}{exo:b2:alkene-redox:8}
Propanone (3 C) and butanal (4 C) join at their carbonyl carbons: \ce{(CH3)2C=CH-CH2CH2CH3},
2-methylhex-2-ene.
\end{solution}

\begin{solution}{exo:b2:alkene-redox:9}
\ce{BH3} (or a bulky borane), then \ce{H2O2} and \ce{NaOH}: hexan-1-ol. Acid-catalysed hydration
goes through the secondary carbocation and gives hexan-2-ol (Markovnikov), with
rearranged by-products.
\end{solution}

\begin{solution}{exo:b2:alkene-redox:10}
The \omterm{def:b2:alkene-redox:epoxide}{peroxy acid}, an electrophile, reacts faster with the more substituted, more electron-rich
ring double bond. The bulky borane reacts with the less hindered terminal \ce{C=CH2}.
\end{solution}

\begin{solution}{exo:b2:alkene-redox:11}
trans: \omterm{def:b2:alkene-redox:epoxide}{peroxy acid}, then water in acid (anti opening). cis: catalytic \ce{OsO4} with a
reoxidant, or cold dilute permanganate (syn).
\end{solution}

\begin{solution}{exo:b2:alkene-redox:12}
\ce{C6H10} has two degrees of unsaturation; one equivalent of \ce{H2} means one \ce{C=C}, so
one ring. A single six-carbon dialdehyde comes from a ring opened at its double bond:
cyclohexene.
\end{solution}

\begin{solution}{pb:b2:alkene-redox:1}
\textbf{1.} $10 \times 12.011 + 16 \times 1.008 = \qty{136.24}{g/mol}$.
\textbf{2.} $(2 \times 10 + 2 - 16)/2 = 3$.
\textbf{3.} One ring and two double bonds.
\textbf{4.} \ce{C10H16 + 2H2 -> C10H20}, 1-isopropyl-4-methylcyclohexane.
\textbf{5.} $1.00/136.24 = \qty{7.34}{mmol}$ of limonene, \qty{14.68}{mmol} of \ce{H2}.
\textbf{6.} $V = nRT/p = 0.01468 \times 8.314 \times 298.15/10^5 = \qty{3.64e-4}{m^3}$,
\qty{364}{mL}.
\textbf{7.} No: the ring carbons 1 and 4 each carry two identical ring paths; the cis and
trans isomers are achiral.
\textbf{8.} Methanal \ce{HCHO} (1 C) from the \ce{=CH2} end, and a single \ce{C9} compound,
3-acetyl-6-oxoheptanal \ce{CH3CO-CH2CH2-CH(COCH3)-CH2-CHO}.
\textbf{9.} Its two carbons belong to the same ring: breaking the double bond opens the ring
without separating anything.
\textbf{10.} Methanal (a gas).
\textbf{11.} The aldehyde groups become carboxylic acids (methanal goes on to \ce{CO2}); the
ketones are unchanged.
\textbf{12.} The ketone and aldehyde of the \ce{C9} chain were joined in the ring; the second
ketone (\ce{COCH3}) and methanal were the two ends of the side chain's double bond.
\textbf{13.} One (C4 of the ring); it is untouched by \omterm{def:b2:alkene-redox:cleavage}{ozonolysis} and survives as a
stereocentre of the \ce{C9} product.
\textbf{14.} The exocyclic \ce{C=CH2}, less hindered than the trisubstituted ring double bond.
\textbf{15.} 2-(4-methylcyclohex-3-en-1-yl)propan-1-ol: \ce{CH2OH} at the former terminal
carbon.
\textbf{16.} Primary.
\textbf{17.} The \ce{OH} goes to the less substituted carbon, opposite to the orientation
predicted by Markovnikov's rule for hydration.
\textbf{18.} Two configurations at the new centre, giving two diastereomers (with the existing
centre).
\textbf{19.} Acid-catalysed hydration (or another Markovnikov hydration): the tertiary alcohol,
$\alpha$-terpineol.
\textbf{20.} The ring double bond: trisubstituted, more electron-rich.
\textbf{21.} 1,2-Epoxy-1-methyl-4-(prop-1-en-2-yl)cyclohexane, the oxygen bridging C1 and C2.
\textbf{22.} At C1, the more substituted carbon (bearing the methyl).
\textbf{23.} 1-Methyl-4-(prop-1-en-2-yl)cyclohexane-1,2-diol, the two \ce{OH} trans (anti
opening).
\textbf{24.} Syn \omterm{def:b2:alkene-redox:dihydroxylation}{dihydroxylation} of the ring double bond (catalytic \ce{OsO4}), which gives the cis
diol --- if the reagent can be made to react with that double bond selectively.
\textbf{25.} \qty{1.00}{g} of limonene absorbs $\boldsymbol{\approx \qty{364}{mL}}$ of
dihydrogen.
\end{solution}
